\mxfront{terminale-e3}
\mxnewpage\phantomsection\label{sujet:3-1}\addcontentsline{toc}{chapter}{Devoir commun · Zéro 1}

\section*{Devoir commun · Zéro 1}

\mxbanner{SUJET ZÉRO · ENTRAÎNEMENT}

\textbf{Durée : 2 h · Note sur 20 · \mxTallied{ex-3-1} exercices indépendants.}\\Calculatrice autorisée. Toute réponse doit être justifiée ; les démarches incomplètes seront prises en compte. Donner les valeurs exactes ; arrondir les probabilités à $10^{-6}$ et les autres approximations à $10^{-4}$, sauf précision différente.

\smallskip\textbf{Nom :}\enspace\dotfill\quad\textbf{Classe :}\enspace\dotfill\quad\textbf{Date :}\enspace\dotfill

\mxtally{ex-3-1}
\subsection*{Exercice 1 · Une évolution discrète \hfill 5 points}On considère la suite $(u_n)$ définie par $u_0=2$ et, pour tout entier naturel $n$, $u_{n+1}=\frac{1}{4}u_n+6$.

\begin{enumerate}

\item Calculer $u_1$ et $u_2$. \hfill\textit{(\num{0,5} pt)}

\item Démontrer par récurrence que, pour tout $n\in\N$, $u_n=8-6\left(\frac{1}{4}\right)^n$. \hfill\textit{(2 pt)}

\item En déduire que, pour tout entier naturel $n$, $2\le u_n<8$. \hfill\textit{(1 pt)}

\item Étudier le signe de $u_{n+1}-u_n$ et en déduire le sens de variation de $(u_n)$. \hfill\textit{(1 pt)}

\item Un élève affirme : « Il existe un entier $n$ tel que $u_n=8$. » Cette affirmation est-elle vraie ? Justifier. \hfill\textit{(\num{0,5} pt)}

\end{enumerate}

\mxtally{ex-3-1}
\subsection*{Exercice 2 · Limites calculées et justifiées \hfill 5 points}Pour tout entier $n\ge1$, on définit $u_n=\frac{3n^2+5n-1}{4n^2+2}$ et $v_n=\sqrt{n^2+5n}-n$.

\begin{enumerate}

\item Déterminer la limite de $(u_n)$ en détaillant la transformation de son expression. \hfill\textit{(1 pt)}

\item Montrer que, pour tout entier $n\ge1$, $v_n=\frac{5}{\sqrt{1+\frac{5}n}+1}$. \hfill\textit{(1 pt)}

\item En déduire la limite de $(v_n)$. \hfill\textit{(1 pt)}

\item On définit $w_n=3n^2-5n$ pour $n\ge1$. Déterminer la limite de $(w_n)$. \hfill\textit{(1 pt)}

\item Un élève affirme que la différence de deux suites tendant vers $+\infty$ tend toujours vers $0$. Cette affirmation est-elle vraie ? Justifier. \hfill\textit{(1 pt)}

\end{enumerate}

\mxnewpage\mxcontinued{Devoir commun · Zéro 1 · suite}

\mxtally{ex-3-1}
\subsection*{Exercice 3 · Limites et asymptotes \hfill 5 points}On considère la fonction $f$ définie par $f(x)=\frac{2x+6}{x-2}$. On note $\mathcal C_f$ sa courbe représentative dans un repère du plan.

\begin{enumerate}

\item Déterminer l'ensemble de définition de $f$. Trouver deux réels $A$ et $B$ tels que, pour tout $x\ne2$, $f(x)=A+\frac B{x-2}$. \hfill\textit{(1 pt)}

\item Déterminer les limites de $f$ en $-\infty$ et en $+\infty$. \hfill\textit{(1 pt)}

\item Déterminer les limites de $f$ en $2^-$ et en $2^+$. \hfill\textit{(1 pt)}

\item Donner les équations des asymptotes de $\mathcal C_f$. Préciser la position de la courbe par rapport à son asymptote horizontale. \hfill\textit{(1 pt)}

\item Indépendamment des questions précédentes, déterminer $\displaystyle\lim_{x\to+\infty}x^{2}e^{-2x}$ en justifiant. \hfill\textit{(1 pt)}

\end{enumerate}

\mxtally{ex-3-1}
\subsection*{Exercice 4 · Codes et sélections \hfill 5 points}Un code est une suite ordonnée de $4$ symboles choisis dans un alphabet qui contient $15$ symboles distincts. Les règles concernant les répétitions sont précisées dans chaque question, indépendamment des autres questions.

\begin{enumerate}

\item Combien de codes peut-on former si les répétitions sont autorisées ? \hfill\textit{(1 pt)}

\item Combien de codes peut-on former si chaque symbole doit être utilisé au plus une fois ? \hfill\textit{(1 pt)}

\item Un symbole de l'alphabet est fixé à l'avance. Les répétitions étant autorisées, combien de codes contiennent ce symbole au moins une fois ? \hfill\textit{(1 pt)}

\item Le premier symbole du code est imposé et les répétitions sont interdites. Combien de codes peut-on former ? \hfill\textit{(1 pt)}

\item Deux symboles distincts de l'alphabet sont interdits. Les autres symboles peuvent être répétés. Combien de codes peut-on former ? \hfill\textit{(1 pt)}

\end{enumerate}

\ifmxteacher\mxnewpage

\section*{Corrigé · Devoir commun · Zéro 1}\label{cor:3-1}

\subsection*{Exercice 1 · Une évolution discrète}\begin{enumerate}

\item \textit{[\num{0,5} pt]} La relation est appliquée successivement aux rangs $0$ et $1$ :\[u_1=\frac{1}{4}\times 2+6=\frac{13}{2},\qquad u_2=\frac{1}{4}\times \frac{13}{2}+6=\frac{61}{8}.\]

\item \textit{[2 pt]} Pour $n\in\N$, notons $P(n)$ la propriété $u_n=8-6\left(\frac{1}{4}\right)^n$.
\par\smallskip\textbf{Initialisation.} Pour $n=0$, $8-6\times1=2=u_0$. Donc $P(0)$ est vraie.
\par\smallskip\textbf{Hérédité.} Fixons un entier $n\ge0$ et supposons $P(n)$ vraie. Par la relation définissant la suite, puis par l'hypothèse de récurrence,
\[\begin{aligned}u_{n+1}&=\frac{1}{4}u_n+6\\
&=\frac{1}{4}\left[8-6\left(\frac{1}{4}\right)^n\right]+6\\
&=2+6-6\left(\frac{1}{4}\right)^{n+1}\\
&=8-6\left(\frac{1}{4}\right)^{n+1}.
\end{aligned}\]
C'est $P(n+1)$ : la propriété est héréditaire.
\par\smallskip\textbf{Conclusion.} Le principe de récurrence donne $u_n=8-6\left(\frac{1}{4}\right)^n$ pour tout $n\in\N$.

\item \textit{[1 pt]} Puisque $0<\frac{1}{4}<1$, pour tout $n\in\N$ on a $0<(\frac{1}{4})^n\le1$. La multiplication par $-6<0$ inverse les inégalités :
\[-6\le-6(\frac{1}{4})^n<0.\]
En ajoutant $8$ et en utilisant la formule précédente, on obtient $2\le u_n<8$.

\item \textit{[1 pt]} Pour tout $n\in\N$, la formule explicite donne
\[\begin{aligned}u_{n+1}-u_n&=-6(\frac{1}{4})^{n+1}+6(\frac{1}{4})^n\\&=6(1-\frac{1}{4})(\frac{1}{4})^n=\frac{9}{2}(\frac{1}{4})^n>0.\end{aligned}\]
La suite est donc strictement croissante.

\item \textit{[\num{0,5} pt]} L'affirmation est fausse. Pour chaque entier $n$, $8-u_n=6(\frac{1}{4})^n>0$. Aucun terme n'est donc égal à $8$, quel que soit le nombre d'étapes.

\end{enumerate}

\subsection*{Exercice 2 · Limites calculées et justifiées}\begin{enumerate}

\item \textit{[1 pt]} Pour $n\ge1$, on divise par $n^2\ne0$ :
\[u_n=\frac{3+\frac{5}{n}-\frac{1}{n^2}}{4+\frac{2}{n^2}}.\]
Les termes $\frac{1}{n}$ et $\frac{1}{n^2}$ tendent vers $0$. Le numérateur tend vers $3$ et le dénominateur vers $4\ne0$ ; par quotient, $u_n\to\frac{3}{4}$.

\item \textit{[1 pt]} Le conjugué est positif. Pour $n\ge1$, $\sqrt{n^2}=n$, d'où
\[\begin{aligned}v_n&=\frac{(n^2+5n)-n^2}{\sqrt{n^2+5n}+n}\\&=\frac{5n}{n(\sqrt{1+\frac{5}{n}}+1)}=\frac{5}{\sqrt{1+\frac{5}{n}}+1}.\end{aligned}\]

\item \textit{[1 pt]} Le dénominateur tend vers $2$ ; limite $\frac{5}{2}$.

\item \textit{[1 pt]} $w_n=n^2(3-\frac{5}{n})\to+\infty$ ; le second facteur tend vers $3>0$.

\item \textit{[1 pt]} Non : $(2n)-n=n\to+\infty$. Le signe $\infty-\infty$ est une forme indéterminée, pas une règle de calcul.

\end{enumerate}

\subsection*{Exercice 3 · Limites et asymptotes}\begin{enumerate}

\item \textit{[1 pt]} Le dénominateur s'annule seulement en $2$ : $D_f=\R\setminus\{2\}$. Comme $2x+6=2(x-2)+10$, on obtient $f(x)=2+\dfrac{10}{x-2}$.

\item \textit{[1 pt]} Les deux limites valent $2$.

\item \textit{[1 pt]} Quand $x\to 2^-$, $x-2\to0$ par valeurs négatives ; comme $10>0$, $\frac{10}{x-2}\to-\infty$, donc $f(x)\to-\infty$. À droite, le dénominateur tend vers $0$ par valeurs positives : le quotient, puis $f(x)$, tendent vers $+\infty$.

\item \textit{[1 pt]} $x=2$ et $y=2$ ; courbe en dessous pour $x<2$ et au-dessus pour $x>2$.

\item \textit{[1 pt]} Posons $t=2x$. Alors $t\to+\infty$ et $x^2e^{-2x}=\frac{1}{4}\frac{t^2}{e^t}$. La croissance comparée donne $\frac{t^2}{e^t}\to0$ ; la limite cherchée vaut donc $0$.

\end{enumerate}

\subsection*{Exercice 4 · Codes et sélections}\begin{enumerate}

\item \textit{[1 pt]} $15^4=50625$ par le principe multiplicatif.

\item \textit{[1 pt]} Les choix successifs ont $15$, $14$, puis un choix de moins à chaque étape, jusqu'à $12$ possibilités. Par multiplication,
\[15\times14\times13\times12=32760.\]
L'ordre des symboles fait partie du code.

\item \textit{[1 pt]} $15^4-14^4=12209$ par complément.

\item \textit{[1 pt]} Le symbole de la première position est fixé et ne peut plus être utilisé. Les $3$ positions suivantes donnent successivement $14,\ldots,12$ choix :
\[14\times13\times12=2184.\]

\item \textit{[1 pt]} $13^4=28561$.

\end{enumerate}

\fi

\mxnewpage\phantomsection\label{sujet:3-2}\addcontentsline{toc}{chapter}{Devoir commun · Zéro 2}

\section*{Devoir commun · Zéro 2}

\mxbanner{SUJET ZÉRO · ENTRAÎNEMENT}

\textbf{Durée : 2 h · Note sur 20 · \mxTallied{ex-3-2} exercices indépendants.}\\Calculatrice autorisée. Toute réponse doit être justifiée ; les démarches incomplètes seront prises en compte. Donner les valeurs exactes ; arrondir les probabilités à $10^{-6}$ et les autres approximations à $10^{-4}$, sauf précision différente.

\smallskip\textbf{Nom :}\enspace\dotfill\quad\textbf{Classe :}\enspace\dotfill\quad\textbf{Date :}\enspace\dotfill

\mxtally{ex-3-2}
\subsection*{Exercice 1 · Une évolution discrète \hfill 5 points}On considère la suite $(u_n)$ définie par $u_0=3$ et, pour tout entier naturel $n$, $u_{n+1}=\frac{1}{3}u_n+\frac{20}{3}$.

\begin{enumerate}

\item Calculer $u_1$ et $u_2$. \hfill\textit{(\num{0,5} pt)}

\item Démontrer par récurrence que, pour tout $n\in\N$, $u_n=10-7\left(\frac{1}{3}\right)^n$. \hfill\textit{(2 pt)}

\item En déduire que, pour tout entier naturel $n$, $3\le u_n<10$. \hfill\textit{(1 pt)}

\item Étudier le signe de $u_{n+1}-u_n$ et en déduire le sens de variation de $(u_n)$. \hfill\textit{(1 pt)}

\item Un élève affirme : « Il existe un entier $n$ tel que $u_n=10$. » Cette affirmation est-elle vraie ? Justifier. \hfill\textit{(\num{0,5} pt)}

\end{enumerate}

\mxtally{ex-3-2}
\subsection*{Exercice 2 · Limites calculées et justifiées \hfill 5 points}Pour tout entier $n\ge1$, on définit $u_n=\frac{2n^2+7n-1}{3n^2+3}$ et $v_n=\sqrt{n^2+7n}-n$.

\begin{enumerate}

\item Déterminer la limite de $(u_n)$ en détaillant la transformation de son expression. \hfill\textit{(1 pt)}

\item Montrer que, pour tout entier $n\ge1$, $v_n=\frac{7}{\sqrt{1+\frac{7}n}+1}$. \hfill\textit{(1 pt)}

\item En déduire la limite de $(v_n)$. \hfill\textit{(1 pt)}

\item On définit $w_n=2n^2-7n$ pour $n\ge1$. Déterminer la limite de $(w_n)$. \hfill\textit{(1 pt)}

\item Un élève affirme que la différence de deux suites tendant vers $+\infty$ tend toujours vers $0$. Cette affirmation est-elle vraie ? Justifier. \hfill\textit{(1 pt)}

\end{enumerate}

\mxnewpage\mxcontinued{Devoir commun · Zéro 2 · suite}

\mxtally{ex-3-2}
\subsection*{Exercice 3 · Limites et asymptotes \hfill 5 points}On considère la fonction $f$ définie par $f(x)=\frac{5x+7}{x-3}$. On note $\mathcal C_f$ sa courbe représentative dans un repère du plan.

\begin{enumerate}

\item Déterminer l'ensemble de définition de $f$. Trouver deux réels $A$ et $B$ tels que, pour tout $x\ne3$, $f(x)=A+\frac B{x-3}$. \hfill\textit{(1 pt)}

\item Déterminer les limites de $f$ en $-\infty$ et en $+\infty$. \hfill\textit{(1 pt)}

\item Déterminer les limites de $f$ en $3^-$ et en $3^+$. \hfill\textit{(1 pt)}

\item Donner les équations des asymptotes de $\mathcal C_f$. Préciser la position de la courbe par rapport à son asymptote horizontale. \hfill\textit{(1 pt)}

\item Indépendamment des questions précédentes, déterminer $\displaystyle\lim_{x\to+\infty}x^{5}e^{-3x}$ en justifiant. \hfill\textit{(1 pt)}

\end{enumerate}

\mxtally{ex-3-2}
\subsection*{Exercice 4 · Codes et sélections \hfill 5 points}Un code est une suite ordonnée de $5$ symboles choisis dans un alphabet qui contient $16$ symboles distincts. Les règles concernant les répétitions sont précisées dans chaque question, indépendamment des autres questions.

\begin{enumerate}

\item Combien de codes peut-on former si les répétitions sont autorisées ? \hfill\textit{(1 pt)}

\item Combien de codes peut-on former si chaque symbole doit être utilisé au plus une fois ? \hfill\textit{(1 pt)}

\item Un symbole de l'alphabet est fixé à l'avance. Les répétitions étant autorisées, combien de codes contiennent ce symbole au moins une fois ? \hfill\textit{(1 pt)}

\item Le premier symbole du code est imposé et les répétitions sont interdites. Combien de codes peut-on former ? \hfill\textit{(1 pt)}

\item Deux symboles distincts de l'alphabet sont interdits. Les autres symboles peuvent être répétés. Combien de codes peut-on former ? \hfill\textit{(1 pt)}

\end{enumerate}

\ifmxteacher\mxnewpage

\section*{Corrigé · Devoir commun · Zéro 2}\label{cor:3-2}

\subsection*{Exercice 1 · Une évolution discrète}\begin{enumerate}

\item \textit{[\num{0,5} pt]} La relation est appliquée successivement aux rangs $0$ et $1$ :\[u_1=\frac{1}{3}\times 3+\frac{20}{3}=\frac{23}{3},\qquad u_2=\frac{1}{3}\times \frac{23}{3}+\frac{20}{3}=\frac{83}{9}.\]

\item \textit{[2 pt]} Pour $n\in\N$, notons $P(n)$ la propriété $u_n=10-7\left(\frac{1}{3}\right)^n$.
\par\smallskip\textbf{Initialisation.} Pour $n=0$, $10-7\times1=3=u_0$. Donc $P(0)$ est vraie.
\par\smallskip\textbf{Hérédité.} Fixons un entier $n\ge0$ et supposons $P(n)$ vraie. Par la relation définissant la suite, puis par l'hypothèse de récurrence,
\[\begin{aligned}u_{n+1}&=\frac{1}{3}u_n+\frac{20}{3}\\
&=\frac{1}{3}\left[10-7\left(\frac{1}{3}\right)^n\right]+\frac{20}{3}\\
&=\frac{10}{3}+\frac{20}{3}-7\left(\frac{1}{3}\right)^{n+1}\\
&=10-7\left(\frac{1}{3}\right)^{n+1}.
\end{aligned}\]
C'est $P(n+1)$ : la propriété est héréditaire.
\par\smallskip\textbf{Conclusion.} Le principe de récurrence donne $u_n=10-7\left(\frac{1}{3}\right)^n$ pour tout $n\in\N$.

\item \textit{[1 pt]} Puisque $0<\frac{1}{3}<1$, pour tout $n\in\N$ on a $0<(\frac{1}{3})^n\le1$. La multiplication par $-7<0$ inverse les inégalités :
\[-7\le-7(\frac{1}{3})^n<0.\]
En ajoutant $10$ et en utilisant la formule précédente, on obtient $3\le u_n<10$.

\item \textit{[1 pt]} Pour tout $n\in\N$, la formule explicite donne
\[\begin{aligned}u_{n+1}-u_n&=-7(\frac{1}{3})^{n+1}+7(\frac{1}{3})^n\\&=7(1-\frac{1}{3})(\frac{1}{3})^n=\frac{14}{3}(\frac{1}{3})^n>0.\end{aligned}\]
La suite est donc strictement croissante.

\item \textit{[\num{0,5} pt]} L'affirmation est fausse. Pour chaque entier $n$, $10-u_n=7(\frac{1}{3})^n>0$. Aucun terme n'est donc égal à $10$, quel que soit le nombre d'étapes.

\end{enumerate}

\subsection*{Exercice 2 · Limites calculées et justifiées}\begin{enumerate}

\item \textit{[1 pt]} Pour $n\ge1$, on divise par $n^2\ne0$ :
\[u_n=\frac{2+\frac{7}{n}-\frac{1}{n^2}}{3+\frac{3}{n^2}}.\]
Les termes $\frac{1}{n}$ et $\frac{1}{n^2}$ tendent vers $0$. Le numérateur tend vers $2$ et le dénominateur vers $3\ne0$ ; par quotient, $u_n\to\frac{2}{3}$.

\item \textit{[1 pt]} Le conjugué est positif. Pour $n\ge1$, $\sqrt{n^2}=n$, d'où
\[\begin{aligned}v_n&=\frac{(n^2+7n)-n^2}{\sqrt{n^2+7n}+n}\\&=\frac{7n}{n(\sqrt{1+\frac{7}{n}}+1)}=\frac{7}{\sqrt{1+\frac{7}{n}}+1}.\end{aligned}\]

\item \textit{[1 pt]} Le dénominateur tend vers $2$ ; limite $\frac{7}{2}$.

\item \textit{[1 pt]} $w_n=n^2(2-\frac{7}{n})\to+\infty$ ; le second facteur tend vers $2>0$.

\item \textit{[1 pt]} Non : $(2n)-n=n\to+\infty$. Le signe $\infty-\infty$ est une forme indéterminée, pas une règle de calcul.

\end{enumerate}

\subsection*{Exercice 3 · Limites et asymptotes}\begin{enumerate}

\item \textit{[1 pt]} Le dénominateur s'annule seulement en $3$ : $D_f=\R\setminus\{3\}$. Comme $5x+7=5(x-3)+22$, on obtient $f(x)=5+\dfrac{22}{x-3}$.

\item \textit{[1 pt]} Les deux limites valent $5$.

\item \textit{[1 pt]} Quand $x\to 3^-$, $x-3\to0$ par valeurs négatives ; comme $22>0$, $\frac{22}{x-3}\to-\infty$, donc $f(x)\to-\infty$. À droite, le dénominateur tend vers $0$ par valeurs positives : le quotient, puis $f(x)$, tendent vers $+\infty$.

\item \textit{[1 pt]} $x=3$ et $y=5$ ; courbe en dessous pour $x<3$ et au-dessus pour $x>3$.

\item \textit{[1 pt]} Posons $t=3x$. Alors $t\to+\infty$ et $x^5e^{-3x}=\frac{1}{243}\frac{t^5}{e^t}$. La croissance comparée donne $\frac{t^5}{e^t}\to0$ ; la limite cherchée vaut donc $0$.

\end{enumerate}

\subsection*{Exercice 4 · Codes et sélections}\begin{enumerate}

\item \textit{[1 pt]} $16^5=1048576$ par le principe multiplicatif.

\item \textit{[1 pt]} Les choix successifs ont $16$, $15$, puis un choix de moins à chaque étape, jusqu'à $12$ possibilités. Par multiplication,
\[16\times15\times14\times13\times12=524160.\]
L'ordre des symboles fait partie du code.

\item \textit{[1 pt]} $16^5-15^5=289201$ par complément.

\item \textit{[1 pt]} Le symbole de la première position est fixé et ne peut plus être utilisé. Les $4$ positions suivantes donnent successivement $15,\ldots,12$ choix :
\[15\times14\times13\times12=32760.\]

\item \textit{[1 pt]} $14^5=537824$.

\end{enumerate}

\fi

\ifmxteacher

\mxnewpage\phantomsection\label{sujet:3-3}\addcontentsline{toc}{chapter}{Devoir commun · A}

\section*{Devoir commun · A}

\mxbanner*{SUJET RÉSERVÉ À L'ÉVALUATION}

\textbf{Durée : 2 h · Note sur 20 · \mxTallied{ex-3-3} exercices indépendants.}\\Calculatrice autorisée. Toute réponse doit être justifiée ; les démarches incomplètes seront prises en compte. Donner les valeurs exactes ; arrondir les probabilités à $10^{-6}$ et les autres approximations à $10^{-4}$, sauf précision différente.

\smallskip\textbf{Nom :}\enspace\dotfill\quad\textbf{Classe :}\enspace\dotfill\quad\textbf{Date :}\enspace\dotfill

\mxtally{ex-3-3}
\subsection*{Exercice 1 · Une évolution discrète \hfill 5 points}On considère la suite $(u_n)$ définie par $u_0=1$ et, pour tout entier naturel $n$, $u_{n+1}=\frac{1}{2}u_n+6$.

\begin{enumerate}

\item Calculer $u_1$ et $u_2$. \hfill\textit{(\num{0,5} pt)}

\item Démontrer par récurrence que, pour tout $n\in\N$, $u_n=12-11\left(\frac{1}{2}\right)^n$. \hfill\textit{(2 pt)}

\item En déduire que, pour tout entier naturel $n$, $1\le u_n<12$. \hfill\textit{(1 pt)}

\item Étudier le signe de $u_{n+1}-u_n$ et en déduire le sens de variation de $(u_n)$. \hfill\textit{(1 pt)}

\item Un élève affirme : « Il existe un entier $n$ tel que $u_n=12$. » Cette affirmation est-elle vraie ? Justifier. \hfill\textit{(\num{0,5} pt)}

\end{enumerate}

\mxtally{ex-3-3}
\subsection*{Exercice 2 · Limites calculées et justifiées \hfill 5 points}Pour tout entier $n\ge1$, on définit $u_n=\frac{5n^2+8n-1}{6n^2+1}$ et $v_n=\sqrt{n^2+8n}-n$.

\begin{enumerate}

\item Déterminer la limite de $(u_n)$ en détaillant la transformation de son expression. \hfill\textit{(1 pt)}

\item Montrer que, pour tout entier $n\ge1$, $v_n=\frac{8}{\sqrt{1+\frac{8}n}+1}$. \hfill\textit{(1 pt)}

\item En déduire la limite de $(v_n)$. \hfill\textit{(1 pt)}

\item On définit $w_n=5n^2-8n$ pour $n\ge1$. Déterminer la limite de $(w_n)$. \hfill\textit{(1 pt)}

\item Un élève affirme que la différence de deux suites tendant vers $+\infty$ tend toujours vers $0$. Cette affirmation est-elle vraie ? Justifier. \hfill\textit{(1 pt)}

\end{enumerate}

\mxnewpage\mxcontinued{Devoir commun · A · suite}

\mxtally{ex-3-3}
\subsection*{Exercice 3 · Limites et asymptotes \hfill 5 points}On considère la fonction $f$ définie par $f(x)=\frac{4x+4}{x-1}$. On note $\mathcal C_f$ sa courbe représentative dans un repère du plan.

\begin{enumerate}

\item Déterminer l'ensemble de définition de $f$. Trouver deux réels $A$ et $B$ tels que, pour tout $x\ne1$, $f(x)=A+\frac B{x-1}$. \hfill\textit{(1 pt)}

\item Déterminer les limites de $f$ en $-\infty$ et en $+\infty$. \hfill\textit{(1 pt)}

\item Déterminer les limites de $f$ en $1^-$ et en $1^+$. \hfill\textit{(1 pt)}

\item Donner les équations des asymptotes de $\mathcal C_f$. Préciser la position de la courbe par rapport à son asymptote horizontale. \hfill\textit{(1 pt)}

\item Indépendamment des questions précédentes, déterminer $\displaystyle\lim_{x\to+\infty}x^{4}e^{-1x}$ en justifiant. \hfill\textit{(1 pt)}

\end{enumerate}

\mxtally{ex-3-3}
\subsection*{Exercice 4 · Codes et sélections \hfill 5 points}Un code est une suite ordonnée de $3$ symboles choisis dans un alphabet qui contient $12$ symboles distincts. Les règles concernant les répétitions sont précisées dans chaque question, indépendamment des autres questions.

\begin{enumerate}

\item Combien de codes peut-on former si les répétitions sont autorisées ? \hfill\textit{(1 pt)}

\item Combien de codes peut-on former si chaque symbole doit être utilisé au plus une fois ? \hfill\textit{(1 pt)}

\item Un symbole de l'alphabet est fixé à l'avance. Les répétitions étant autorisées, combien de codes contiennent ce symbole au moins une fois ? \hfill\textit{(1 pt)}

\item Le premier symbole du code est imposé et les répétitions sont interdites. Combien de codes peut-on former ? \hfill\textit{(1 pt)}

\item Deux symboles distincts de l'alphabet sont interdits. Les autres symboles peuvent être répétés. Combien de codes peut-on former ? \hfill\textit{(1 pt)}

\end{enumerate}

\ifmxteacher\mxnewpage

\section*{Corrigé · Devoir commun · A}\label{cor:3-3}

\subsection*{Exercice 1 · Une évolution discrète}\begin{enumerate}

\item \textit{[\num{0,5} pt]} La relation est appliquée successivement aux rangs $0$ et $1$ :\[u_1=\frac{1}{2}\times 1+6=\frac{13}{2},\qquad u_2=\frac{1}{2}\times \frac{13}{2}+6=\frac{37}{4}.\]

\item \textit{[2 pt]} Pour $n\in\N$, notons $P(n)$ la propriété $u_n=12-11\left(\frac{1}{2}\right)^n$.
\par\smallskip\textbf{Initialisation.} Pour $n=0$, $12-11\times1=1=u_0$. Donc $P(0)$ est vraie.
\par\smallskip\textbf{Hérédité.} Fixons un entier $n\ge0$ et supposons $P(n)$ vraie. Par la relation définissant la suite, puis par l'hypothèse de récurrence,
\[\begin{aligned}u_{n+1}&=\frac{1}{2}u_n+6\\
&=\frac{1}{2}\left[12-11\left(\frac{1}{2}\right)^n\right]+6\\
&=6+6-11\left(\frac{1}{2}\right)^{n+1}\\
&=12-11\left(\frac{1}{2}\right)^{n+1}.
\end{aligned}\]
C'est $P(n+1)$ : la propriété est héréditaire.
\par\smallskip\textbf{Conclusion.} Le principe de récurrence donne $u_n=12-11\left(\frac{1}{2}\right)^n$ pour tout $n\in\N$.

\item \textit{[1 pt]} Puisque $0<\frac{1}{2}<1$, pour tout $n\in\N$ on a $0<(\frac{1}{2})^n\le1$. La multiplication par $-11<0$ inverse les inégalités :
\[-11\le-11(\frac{1}{2})^n<0.\]
En ajoutant $12$ et en utilisant la formule précédente, on obtient $1\le u_n<12$.

\item \textit{[1 pt]} Pour tout $n\in\N$, la formule explicite donne
\[\begin{aligned}u_{n+1}-u_n&=-11(\frac{1}{2})^{n+1}+11(\frac{1}{2})^n\\&=11(1-\frac{1}{2})(\frac{1}{2})^n=\frac{11}{2}(\frac{1}{2})^n>0.\end{aligned}\]
La suite est donc strictement croissante.

\item \textit{[\num{0,5} pt]} L'affirmation est fausse. Pour chaque entier $n$, $12-u_n=11(\frac{1}{2})^n>0$. Aucun terme n'est donc égal à $12$, quel que soit le nombre d'étapes.

\end{enumerate}

\subsection*{Exercice 2 · Limites calculées et justifiées}\begin{enumerate}

\item \textit{[1 pt]} Pour $n\ge1$, on divise par $n^2\ne0$ :
\[u_n=\frac{5+\frac{8}{n}-\frac{1}{n^2}}{6+\frac{1}{n^2}}.\]
Les termes $\frac{1}{n}$ et $\frac{1}{n^2}$ tendent vers $0$. Le numérateur tend vers $5$ et le dénominateur vers $6\ne0$ ; par quotient, $u_n\to\frac{5}{6}$.

\item \textit{[1 pt]} Le conjugué est positif. Pour $n\ge1$, $\sqrt{n^2}=n$, d'où
\[\begin{aligned}v_n&=\frac{(n^2+8n)-n^2}{\sqrt{n^2+8n}+n}\\&=\frac{8n}{n(\sqrt{1+\frac{8}{n}}+1)}=\frac{8}{\sqrt{1+\frac{8}{n}}+1}.\end{aligned}\]

\item \textit{[1 pt]} Le dénominateur tend vers $2$ ; limite $4$.

\item \textit{[1 pt]} $w_n=n^2(5-\frac{8}{n})\to+\infty$ ; le second facteur tend vers $5>0$.

\item \textit{[1 pt]} Non : $(2n)-n=n\to+\infty$. Le signe $\infty-\infty$ est une forme indéterminée, pas une règle de calcul.

\end{enumerate}

\subsection*{Exercice 3 · Limites et asymptotes}\begin{enumerate}

\item \textit{[1 pt]} Le dénominateur s'annule seulement en $1$ : $D_f=\R\setminus\{1\}$. Comme $4x+4=4(x-1)+8$, on obtient $f(x)=4+\dfrac{8}{x-1}$.

\item \textit{[1 pt]} Les deux limites valent $4$.

\item \textit{[1 pt]} Quand $x\to 1^-$, $x-1\to0$ par valeurs négatives ; comme $8>0$, $\frac{8}{x-1}\to-\infty$, donc $f(x)\to-\infty$. À droite, le dénominateur tend vers $0$ par valeurs positives : le quotient, puis $f(x)$, tendent vers $+\infty$.

\item \textit{[1 pt]} $x=1$ et $y=4$ ; courbe en dessous pour $x<1$ et au-dessus pour $x>1$.

\item \textit{[1 pt]} La croissance comparée de la puissance et de l'exponentielle donne directement $x^4e^{-x}=\frac{x^4}{e^x}\to0$ lorsque $x\to+\infty$.

\end{enumerate}

\subsection*{Exercice 4 · Codes et sélections}\begin{enumerate}

\item \textit{[1 pt]} $12^3=1728$ par le principe multiplicatif.

\item \textit{[1 pt]} Les choix successifs ont $12$, $11$, puis un choix de moins à chaque étape, jusqu'à $10$ possibilités. Par multiplication,
\[12\times11\times10=1320.\]
L'ordre des symboles fait partie du code.

\item \textit{[1 pt]} $12^3-11^3=397$ par complément.

\item \textit{[1 pt]} Le symbole de la première position est fixé et ne peut plus être utilisé. Les $2$ positions suivantes donnent successivement $11,\ldots,10$ choix :
\[11\times10=110.\]

\item \textit{[1 pt]} $10^3=1000$.

\end{enumerate}

\fi

\fi

\ifmxteacher

\mxnewpage\phantomsection\label{sujet:3-4}\addcontentsline{toc}{chapter}{Devoir commun · B}

\section*{Devoir commun · B}

\mxbanner*{SUJET RÉSERVÉ À L'ÉVALUATION}

\textbf{Durée : 2 h · Note sur 20 · \mxTallied{ex-3-4} exercices indépendants.}\\Calculatrice autorisée. Toute réponse doit être justifiée ; les démarches incomplètes seront prises en compte. Donner les valeurs exactes ; arrondir les probabilités à $10^{-6}$ et les autres approximations à $10^{-4}$, sauf précision différente.

\smallskip\textbf{Nom :}\enspace\dotfill\quad\textbf{Classe :}\enspace\dotfill\quad\textbf{Date :}\enspace\dotfill

\mxtally{ex-3-4}
\subsection*{Exercice 1 · Une évolution discrète \hfill 5 points}On considère la suite $(u_n)$ définie par $u_0=2$ et, pour tout entier naturel $n$, $u_{n+1}=\frac{1}{5}u_n+\frac{32}{5}$.

\begin{enumerate}

\item Calculer $u_1$ et $u_2$. \hfill\textit{(\num{0,5} pt)}

\item Démontrer par récurrence que, pour tout $n\in\N$, $u_n=8-6\left(\frac{1}{5}\right)^n$. \hfill\textit{(2 pt)}

\item En déduire que, pour tout entier naturel $n$, $2\le u_n<8$. \hfill\textit{(1 pt)}

\item Étudier le signe de $u_{n+1}-u_n$ et en déduire le sens de variation de $(u_n)$. \hfill\textit{(1 pt)}

\item Un élève affirme : « Il existe un entier $n$ tel que $u_n=8$. » Cette affirmation est-elle vraie ? Justifier. \hfill\textit{(\num{0,5} pt)}

\end{enumerate}

\mxtally{ex-3-4}
\subsection*{Exercice 2 · Limites calculées et justifiées \hfill 5 points}Pour tout entier $n\ge1$, on définit $u_n=\frac{4n^2+5n-1}{5n^2+2}$ et $v_n=\sqrt{n^2+5n}-n$.

\begin{enumerate}

\item Déterminer la limite de $(u_n)$ en détaillant la transformation de son expression. \hfill\textit{(1 pt)}

\item Montrer que, pour tout entier $n\ge1$, $v_n=\frac{5}{\sqrt{1+\frac{5}n}+1}$. \hfill\textit{(1 pt)}

\item En déduire la limite de $(v_n)$. \hfill\textit{(1 pt)}

\item On définit $w_n=4n^2-5n$ pour $n\ge1$. Déterminer la limite de $(w_n)$. \hfill\textit{(1 pt)}

\item Un élève affirme que la différence de deux suites tendant vers $+\infty$ tend toujours vers $0$. Cette affirmation est-elle vraie ? Justifier. \hfill\textit{(1 pt)}

\end{enumerate}

\mxnewpage\mxcontinued{Devoir commun · B · suite}

\mxtally{ex-3-4}
\subsection*{Exercice 3 · Limites et asymptotes \hfill 5 points}On considère la fonction $f$ définie par $f(x)=\frac{3x+6}{x-2}$. On note $\mathcal C_f$ sa courbe représentative dans un repère du plan.

\begin{enumerate}

\item Déterminer l'ensemble de définition de $f$. Trouver deux réels $A$ et $B$ tels que, pour tout $x\ne2$, $f(x)=A+\frac B{x-2}$. \hfill\textit{(1 pt)}

\item Déterminer les limites de $f$ en $-\infty$ et en $+\infty$. \hfill\textit{(1 pt)}

\item Déterminer les limites de $f$ en $2^-$ et en $2^+$. \hfill\textit{(1 pt)}

\item Donner les équations des asymptotes de $\mathcal C_f$. Préciser la position de la courbe par rapport à son asymptote horizontale. \hfill\textit{(1 pt)}

\item Indépendamment des questions précédentes, déterminer $\displaystyle\lim_{x\to+\infty}x^{3}e^{-2x}$ en justifiant. \hfill\textit{(1 pt)}

\end{enumerate}

\mxtally{ex-3-4}
\subsection*{Exercice 4 · Codes et sélections \hfill 5 points}Un code est une suite ordonnée de $4$ symboles choisis dans un alphabet qui contient $13$ symboles distincts. Les règles concernant les répétitions sont précisées dans chaque question, indépendamment des autres questions.

\begin{enumerate}

\item Combien de codes peut-on former si les répétitions sont autorisées ? \hfill\textit{(1 pt)}

\item Combien de codes peut-on former si chaque symbole doit être utilisé au plus une fois ? \hfill\textit{(1 pt)}

\item Un symbole de l'alphabet est fixé à l'avance. Les répétitions étant autorisées, combien de codes contiennent ce symbole au moins une fois ? \hfill\textit{(1 pt)}

\item Le premier symbole du code est imposé et les répétitions sont interdites. Combien de codes peut-on former ? \hfill\textit{(1 pt)}

\item Deux symboles distincts de l'alphabet sont interdits. Les autres symboles peuvent être répétés. Combien de codes peut-on former ? \hfill\textit{(1 pt)}

\end{enumerate}

\ifmxteacher\mxnewpage

\section*{Corrigé · Devoir commun · B}\label{cor:3-4}

\subsection*{Exercice 1 · Une évolution discrète}\begin{enumerate}

\item \textit{[\num{0,5} pt]} La relation est appliquée successivement aux rangs $0$ et $1$ :\[u_1=\frac{1}{5}\times 2+\frac{32}{5}=\frac{34}{5},\qquad u_2=\frac{1}{5}\times \frac{34}{5}+\frac{32}{5}=\frac{194}{25}.\]

\item \textit{[2 pt]} Pour $n\in\N$, notons $P(n)$ la propriété $u_n=8-6\left(\frac{1}{5}\right)^n$.
\par\smallskip\textbf{Initialisation.} Pour $n=0$, $8-6\times1=2=u_0$. Donc $P(0)$ est vraie.
\par\smallskip\textbf{Hérédité.} Fixons un entier $n\ge0$ et supposons $P(n)$ vraie. Par la relation définissant la suite, puis par l'hypothèse de récurrence,
\[\begin{aligned}u_{n+1}&=\frac{1}{5}u_n+\frac{32}{5}\\
&=\frac{1}{5}\left[8-6\left(\frac{1}{5}\right)^n\right]+\frac{32}{5}\\
&=\frac{8}{5}+\frac{32}{5}-6\left(\frac{1}{5}\right)^{n+1}\\
&=8-6\left(\frac{1}{5}\right)^{n+1}.
\end{aligned}\]
C'est $P(n+1)$ : la propriété est héréditaire.
\par\smallskip\textbf{Conclusion.} Le principe de récurrence donne $u_n=8-6\left(\frac{1}{5}\right)^n$ pour tout $n\in\N$.

\item \textit{[1 pt]} Puisque $0<\frac{1}{5}<1$, pour tout $n\in\N$ on a $0<(\frac{1}{5})^n\le1$. La multiplication par $-6<0$ inverse les inégalités :
\[-6\le-6(\frac{1}{5})^n<0.\]
En ajoutant $8$ et en utilisant la formule précédente, on obtient $2\le u_n<8$.

\item \textit{[1 pt]} Pour tout $n\in\N$, la formule explicite donne
\[\begin{aligned}u_{n+1}-u_n&=-6(\frac{1}{5})^{n+1}+6(\frac{1}{5})^n\\&=6(1-\frac{1}{5})(\frac{1}{5})^n=\frac{24}{5}(\frac{1}{5})^n>0.\end{aligned}\]
La suite est donc strictement croissante.

\item \textit{[\num{0,5} pt]} L'affirmation est fausse. Pour chaque entier $n$, $8-u_n=6(\frac{1}{5})^n>0$. Aucun terme n'est donc égal à $8$, quel que soit le nombre d'étapes.

\end{enumerate}

\subsection*{Exercice 2 · Limites calculées et justifiées}\begin{enumerate}

\item \textit{[1 pt]} Pour $n\ge1$, on divise par $n^2\ne0$ :
\[u_n=\frac{4+\frac{5}{n}-\frac{1}{n^2}}{5+\frac{2}{n^2}}.\]
Les termes $\frac{1}{n}$ et $\frac{1}{n^2}$ tendent vers $0$. Le numérateur tend vers $4$ et le dénominateur vers $5\ne0$ ; par quotient, $u_n\to\frac{4}{5}$.

\item \textit{[1 pt]} Le conjugué est positif. Pour $n\ge1$, $\sqrt{n^2}=n$, d'où
\[\begin{aligned}v_n&=\frac{(n^2+5n)-n^2}{\sqrt{n^2+5n}+n}\\&=\frac{5n}{n(\sqrt{1+\frac{5}{n}}+1)}=\frac{5}{\sqrt{1+\frac{5}{n}}+1}.\end{aligned}\]

\item \textit{[1 pt]} Le dénominateur tend vers $2$ ; limite $\frac{5}{2}$.

\item \textit{[1 pt]} $w_n=n^2(4-\frac{5}{n})\to+\infty$ ; le second facteur tend vers $4>0$.

\item \textit{[1 pt]} Non : $(2n)-n=n\to+\infty$. Le signe $\infty-\infty$ est une forme indéterminée, pas une règle de calcul.

\end{enumerate}

\subsection*{Exercice 3 · Limites et asymptotes}\begin{enumerate}

\item \textit{[1 pt]} Le dénominateur s'annule seulement en $2$ : $D_f=\R\setminus\{2\}$. Comme $3x+6=3(x-2)+12$, on obtient $f(x)=3+\dfrac{12}{x-2}$.

\item \textit{[1 pt]} Les deux limites valent $3$.

\item \textit{[1 pt]} Quand $x\to 2^-$, $x-2\to0$ par valeurs négatives ; comme $12>0$, $\frac{12}{x-2}\to-\infty$, donc $f(x)\to-\infty$. À droite, le dénominateur tend vers $0$ par valeurs positives : le quotient, puis $f(x)$, tendent vers $+\infty$.

\item \textit{[1 pt]} $x=2$ et $y=3$ ; courbe en dessous pour $x<2$ et au-dessus pour $x>2$.

\item \textit{[1 pt]} Posons $t=2x$. Alors $t\to+\infty$ et $x^3e^{-2x}=\frac{1}{8}\frac{t^3}{e^t}$. La croissance comparée donne $\frac{t^3}{e^t}\to0$ ; la limite cherchée vaut donc $0$.

\end{enumerate}

\subsection*{Exercice 4 · Codes et sélections}\begin{enumerate}

\item \textit{[1 pt]} $13^4=28561$ par le principe multiplicatif.

\item \textit{[1 pt]} Les choix successifs ont $13$, $12$, puis un choix de moins à chaque étape, jusqu'à $10$ possibilités. Par multiplication,
\[13\times12\times11\times10=17160.\]
L'ordre des symboles fait partie du code.

\item \textit{[1 pt]} $13^4-12^4=7825$ par complément.

\item \textit{[1 pt]} Le symbole de la première position est fixé et ne peut plus être utilisé. Les $3$ positions suivantes donnent successivement $12,\ldots,10$ choix :
\[12\times11\times10=1320.\]

\item \textit{[1 pt]} $11^4=14641$.

\end{enumerate}

\fi

\fi

